\documentclass[11pt,letterpaper]{article}

% -------------------- PAQUETES --------------------
\usepackage[spanish]{babel}
\usepackage[utf8]{inputenc}
\usepackage[T1]{fontenc}
\usepackage{amsmath,amssymb}
\usepackage{graphicx}
\usepackage{float}
\usepackage[margin=2.5cm]{geometry}
\usepackage{booktabs}
\usepackage{array}
\usepackage{xcolor}
\usepackage{listings}
\usepackage{caption}
\usepackage{microtype}
\usepackage{parskip}
\usepackage{pdfpages}   % para importar la carátula desde un PDF
\usepackage{hyperref}

% -------------------- CONFIGURACIÓN --------------------
\hypersetup{
    colorlinks=true,
    linkcolor=blue,
    urlcolor=blue,
    citecolor=blue
}

\addto\captionsspanish{%
    \renewcommand{\tablename}{Tabla}%
    \renewcommand{\listtablename}{Índice de tablas}%
}
\renewcommand{\lstlistingname}{Código}

% Macro usada en la sección de LeetCode
\newcommand{\bigO}[1]{\ensuremath{\mathcal{O}(#1)}}

\definecolor{codegreen}{rgb}{0,0.5,0}
\definecolor{codegray}{rgb}{0.45,0.45,0.45}
\definecolor{codepurple}{rgb}{0.58,0,0.82}
\definecolor{backcolour}{rgb}{0.96,0.96,0.96}

\lstdefinelanguage{TypeScript}{
    keywords={const,let,function,class,type,if,else,for,return,new,private,void,number,undefined,export,this},
    sensitive=true,
    comment=[l]{//},
    morecomment=[s]{/*}{*/},
    morestring=[b]",
    morestring=[b]'
}

\lstdefinestyle{codigo}{
    backgroundcolor=\color{backcolour},
    commentstyle=\color{codegreen},
    keywordstyle=\color{blue},
    numberstyle=\tiny\color{codegray},
    stringstyle=\color{codepurple},
    basicstyle=\ttfamily\footnotesize,
    breaklines=true,
    captionpos=b,
    keepspaces=true,
    numbers=left,
    numbersep=8pt,
    showspaces=false,
    showstringspaces=false,
    showtabs=false,
    tabsize=4,
    language=TypeScript
}

\lstset{style=codigo}

% -------------------- DOCUMENTO --------------------
\begin{document}

% Carátula: primera página importada desde un PDF
% (cambia "caratula.pdf" por el nombre de tu archivo)
\includepdf[pages=1,fitpaper=true,pagecommand={\thispagestyle{empty}}]{/home/bizueto/documents/caratula_practicas/caratula.pdf}

\section{Comparación simple}

\subsection{Bubble Sort (ordenamiento de burbuja)}
\subsubsection{Funcionamiento}

Bubble Sort recorre el arreglo de izquierda a derecha comparando \textbf{pares de elementos vecinos}. Si el elemento de la izquierda es mayor que el de la derecha, ambos se intercambian. De esta manera, en cada pasada el mayor de los elementos que aún no están acomodados ``sube'' como una burbuja hasta el extremo derecho.

Para un arreglo \texttt{arr} de $n$ elementos, el algoritmo realiza lo siguiente:

\begin{enumerate}
  \item Se repiten $n-1$ pasadas, indexadas por $i = 0, 1, \dots, n-2$.
  \item En la pasada $i$, el índice $j$ recorre las posiciones $0$ a $n-i-2$. Para cada $j$ se compara \texttt{arr[j]} con \texttt{arr[j+1]}.
  \item Si \texttt{arr[j] > arr[j+1]}, los dos elementos se intercambian.
  \item Al terminar la pasada $i$, el elemento más grande de la zona revisada queda en su posición definitiva, por lo que las $i+1$ posiciones finales ya están ordenadas y la siguiente pasada puede ignorarlas.
\end{enumerate}

En el visualizador, cada comparación y cada intercambio se registran como pasos independientes, y los dos vecinos que se comparan se resaltan. Como solo se intercambian elementos cuando \texttt{arr[j] > arr[j+1]} (comparación estricta), los elementos iguales conservan su orden relativo y el algoritmo es \textbf{estable}. Además trabaja \textbf{in situ}, pues los intercambios se hacen dentro del mismo arreglo.

\subsubsection{Complejidad en tiempo}

En la implementación analizada, la pasada $i$ siempre realiza $n-i-1$ comparaciones, sin importar el contenido del arreglo (no existe una bandera que detenga el algoritmo anticipadamente). El número total de comparaciones es, por lo tanto,
\[
  \sum_{i=0}^{n-2} (n-i-1) = \frac{n(n-1)}{2} \in \Theta(n^2).
\]

\paragraph{Mejor caso.} Ocurre cuando el arreglo ya está ordenado. No se realiza ningún intercambio, pero las $\frac{n(n-1)}{2}$ comparaciones se ejecutan de todos modos, por lo que esta versión sigue siendo $\Theta(n^2)$. Si se agrega una bandera que detecte una pasada sin intercambios y termine el ciclo, el mejor caso baja a $\Theta(n)$ (una sola pasada de $n-1$ comparaciones).

\paragraph{Peor caso.} Ocurre cuando el arreglo está en orden descendente. Cada comparación provoca un intercambio, de modo que se realizan $\frac{n(n-1)}{2}$ comparaciones y $\frac{n(n-1)}{2}$ intercambios:
\[
  T_{\text{peor}}(n) \in \Theta(n^2).
\]

\paragraph{Caso promedio.} Para datos aleatorios se hacen $\frac{n(n-1)}{2}$ comparaciones y alrededor de $\frac{n(n-1)}{4}$ intercambios, es decir, $\Theta(n^2)$.

\subsubsection{Complejidad en espacio}

Solo se emplean los índices $i$ y $j$ y una variable temporal para el intercambio, por lo que
\[
  S(n) \in O(1).
\]

\subsection{Insertion Sort (ordenamiento por inserción)}
\subsubsection{Funcionamiento}

El ordenamiento por inserción construye el arreglo ordenado \emph{de izquierda a derecha}, de manera similar a como se ordenan las cartas de una mano. Al inicio, el primer elemento ya se considera una lista ordenada de tamaño uno; a partir de ahí, en cada pasada se toma el siguiente elemento, llamado \textbf{clave}, y se inserta en la posición correcta dentro de la parte que ya está ordenada.

Para un arreglo \texttt{arr} de $n$ elementos, el algoritmo realiza lo siguiente:

\begin{enumerate}
  \item Para cada posición $i$ desde $1$ hasta $n-1$, se guarda la clave \texttt{clave = arr[i]} y se inicializa $j = i - 1$. En ese momento la posición $i$ queda como un ``hueco''.
  \item Mientras $j \geq 0$ y \texttt{arr[j] > clave}, el elemento \texttt{arr[j]} se desplaza una posición a la derecha (\texttt{arr[j+1] = arr[j]}) y se decrementa $j$. De este modo el hueco se va recorriendo hacia la izquierda.
  \item Cuando se encuentra un elemento que no es mayor que la clave, o se llega al inicio del arreglo, se coloca la clave en el hueco: \texttt{arr[j+1] = clave}.
  \item Al terminar la pasada $i$, las posiciones $0$ a $i$ contienen los mismos elementos en orden creciente.
\end{enumerate}

Al concluir la pasada $n-1$ todo el arreglo queda ordenado. En la implementación en TypeScript del visualizador, cada una de estas acciones (tomar la clave, comparar, desplazar e insertar) se registra como un paso independiente, de modo que se puede observar cómo la parte ordenada crece de izquierda a derecha mientras la clave permanece ``levantada'' sobre el hueco.

Como el desplazamiento solo ocurre cuando \texttt{arr[j] > clave} (comparación estricta), los elementos iguales nunca se intercambian entre sí y conservan su orden relativo, por lo que el algoritmo es \textbf{estable}. Además, trabaja \textbf{in situ}: solo requiere una variable auxiliar para la clave.

\subsubsection{Complejidad en tiempo}

El trabajo del algoritmo depende de cuántas veces se ejecuta el ciclo interno (comparaciones y desplazamientos) en cada pasada, lo cual depende a su vez del orden inicial de los datos.

\paragraph{Mejor caso.} Ocurre cuando el arreglo ya está ordenado de forma ascendente. En cada pasada la primera comparación (\texttt{arr[i-1] > clave}) resulta falsa, por lo que el ciclo interno termina de inmediato y no hay desplazamientos. Se realizan $n-1$ comparaciones en total, así que el tiempo es lineal:
\[
  T_{\text{mejor}}(n) = n - 1 \in \Theta(n).
\]

\paragraph{Peor caso.} Ocurre cuando el arreglo está ordenado de forma descendente. En la pasada $i$, la clave es menor que todos los elementos de la parte ordenada, de modo que se desplazan los $i$ elementos y se llega hasta el inicio del arreglo. El número total de comparaciones (y de desplazamientos) es
\[
  T_{\text{peor}}(n) = \sum_{i=1}^{n-1} i = \frac{n(n-1)}{2} \in \Theta(n^2).
\]

\paragraph{Caso promedio.} Para datos en orden aleatorio, en cada pasada la clave se desplaza en promedio la mitad de la parte ordenada, por lo que se hacen aproximadamente $n^2/4$ desplazamientos. El tiempo promedio sigue siendo $\Theta(n^2)$.

\subsubsection{Complejidad en espacio}

El algoritmo solo utiliza un número constante de variables adicionales (\texttt{clave}, $i$ y $j$), sin importar el tamaño de la entrada. Por lo tanto, su complejidad espacial auxiliar es
\[
  S(n) \in O(1).
\]

\subsection{Selection Sort (ordenamiento por selección)}
\subsubsection{Funcionamiento}

Selection Sort divide el arreglo en una parte ordenada, que crece desde la izquierda, y una parte sin ordenar. En cada pasada \textbf{selecciona el elemento mínimo} de la parte sin ordenar y lo coloca en la primera posición de esa parte.

\begin{enumerate}
  \item Para cada posición \texttt{inicio} desde $0$ hasta $n-2$, se supone que el mínimo está en \texttt{inicio} (\texttt{minimo = inicio}).
  \item Se recorren las posiciones \texttt{comparando} desde \texttt{inicio + 1} hasta $n-1$. Si \texttt{lista[comparando] < lista[minimo]}, se actualiza \texttt{minimo = comparando}.
  \item Al terminar el recorrido, se intercambian \texttt{lista[inicio]} y \texttt{lista[minimo]}.
  \item Las posiciones $0$ a \texttt{inicio} quedan ordenadas y contienen sus valores definitivos.
\end{enumerate}

Una característica importante es que cada pasada realiza \textbf{exactamente un intercambio}, a diferencia de Bubble Sort. Sin embargo, ese intercambio puede mover un elemento a una posición lejana y saltarse a otros elementos iguales, por lo que el algoritmo \textbf{no es estable}. Es un algoritmo \textbf{in situ}.

\subsubsection{Complejidad en tiempo}

Para encontrar el mínimo de la pasada con $\texttt{inicio} = i$ siempre se recorren todos los elementos restantes, es decir, se hacen $n-i-1$ comparaciones, independientemente de cómo estén ordenados los datos:
\[
  \sum_{i=0}^{n-2} (n-i-1) = \frac{n(n-1)}{2} \in \Theta(n^2).
\]

\paragraph{Mejor caso.} Con el arreglo ya ordenado, el mínimo de cada pasada ya está en \texttt{inicio}, pero igualmente se hacen todas las comparaciones: $\Theta(n^2)$.

\paragraph{Peor caso.} También $\Theta(n^2)$ comparaciones. Los intercambios son a lo más $n-1$ (uno por pasada), de modo que el costo de movimientos es $O(n)$.

\paragraph{Caso promedio.} $\Theta(n^2)$. Este algoritmo no se beneficia del orden inicial de los datos; su ventaja es el número reducido de intercambios, útil cuando escribir en memoria es costoso.

\subsubsection{Complejidad en espacio}

Se requieren únicamente las variables \texttt{inicio}, \texttt{minimo} y \texttt{comparando}:
\[
  S(n) \in O(1).
\]

\subsection{Comparación entre los algoritmos simples}

\begin{table}[htbp]
  \centering
  \caption{Comparación entre Bubble Sort, Insertion Sort y Selection Sort.}
  \begin{tabular}{lccc}
    \toprule
    \textbf{Característica} & \textbf{Bubble} & \textbf{Insertion} & \textbf{Selection} \\
    \midrule
    Mejor caso       & $\Theta(n^2)^{*}$ & $\Theta(n)$   & $\Theta(n^2)$ \\
    Caso promedio    & $\Theta(n^2)$     & $\Theta(n^2)$ & $\Theta(n^2)$ \\
    Peor caso        & $\Theta(n^2)$     & $\Theta(n^2)$ & $\Theta(n^2)$ \\
    Espacio auxiliar & $O(1)$            & $O(1)$        & $O(1)$ \\
    Estable          & Sí                & Sí            & No \\
    In situ          & Sí                & Sí            & Sí \\
    Estrategia       & Intercambio de vecinos & Inserción & Selección del mínimo \\
    \bottomrule
  \end{tabular}

  \smallskip
  {\footnotesize $^{*}$ $\Theta(n)$ si se añade una bandera que detenga el algoritmo cuando una pasada no hace intercambios.}
\end{table}

Los tres algoritmos tienen tiempo cuadrático en el caso promedio y en el peor caso, y los tres trabajan in situ con memoria adicional constante. Insertion Sort es el único que aprovecha que los datos ya estén casi ordenados, pues en el mejor caso baja a $\Theta(n)$. Bubble Sort, en la versión analizada, siempre realiza todas las comparaciones, y Selection Sort realiza muy pocos intercambios (a lo más $n-1$), lo que resulta útil cuando escribir en memoria es costoso, aunque no es estable.

\clearpage

\section{Comparación eficiente}

\subsection{Merge Sort (ordenamiento por mezcla)}
\subsubsection{Funcionamiento}

Merge Sort sigue la estrategia \textbf{divide y vencerás}, que consta de tres etapas:

\begin{enumerate}
  \item \textbf{Dividir:} el subarreglo \texttt{arr[lo..hi]} se parte en dos mitades a partir de la posición intermedia $\texttt{medio} = \lfloor (lo + hi)/2 \rfloor$. La mitad izquierda abarca de \texttt{lo} a \texttt{medio} y la derecha de \texttt{medio + 1} a \texttt{hi}.
  \item \textbf{Vencer:} cada mitad se ordena de forma recursiva con el mismo procedimiento. El \textbf{caso base} se da cuando \texttt{lo = hi}, es decir, cuando el subarreglo tiene un solo elemento, el cual ya está ordenado por definición.
  \item \textbf{Combinar (mezclar):} las dos mitades, que ya están ordenadas, se fusionan en un solo bloque ordenado.
\end{enumerate}

La operación de mezcla es el núcleo del algoritmo. Se utiliza una lista auxiliar y tres índices: $i$ recorre la mitad izquierda, $j$ la mitad derecha y $k$ la posición de escritura en el auxiliar.

\begin{enumerate}
  \item Mientras queden elementos en ambas mitades, se comparan \texttt{arr[i]} y \texttt{arr[j]}. El menor de los dos se copia al auxiliar y se avanza el índice de la mitad de la que provino (en caso de empate se toma el de la izquierda, pues la condición es $\leq$).
  \item Cuando una de las mitades se agota, los elementos restantes de la otra se copian directamente al auxiliar, ya que están ordenados y son mayores o iguales que todo lo colocado hasta ese momento.
  \item Finalmente, el contenido del auxiliar se copia de regreso a \texttt{arr[lo..hi]}, que queda ordenado.
\end{enumerate}

El procedimiento general consiste entonces en dividir el arreglo hasta obtener elementos individuales y después mezclar bloques cada vez más grandes hasta reconstruir el arreglo completo. Esto es lo que muestra el visualizador: primero las divisiones sucesivas por niveles y luego las mezclas, con la lista auxiliar dibujada debajo del arreglo.

Debido a que en la mezcla se prefiere el elemento de la izquierda en caso de igualdad, Merge Sort es un algoritmo \textbf{estable}.

\subsubsection{Complejidad en tiempo}

La división del arreglo en mitades genera un árbol de recursión de aproximadamente $\lceil \log_2 n \rceil$ niveles. En cada nivel, la mezcla recorre cada elemento una sola vez, lo que representa un trabajo de $\Theta(n)$ por nivel. El tiempo satisface la recurrencia
\[
  T(n) = 2\,T\!\left(\frac{n}{2}\right) + \Theta(n), \qquad T(1) = \Theta(1),
\]
cuya solución, por el Teorema Maestro, es $T(n) = \Theta(n \log n)$.

\paragraph{Mejor caso.} Si el arreglo ya está ordenado, en cada mezcla todos los elementos de la mitad izquierda son menores o iguales que los de la derecha, por lo que solo se hacen alrededor de $\frac{n}{2}\log_2 n$ comparaciones en total. Aun así, todos los elementos se siguen copiando al auxiliar y de regreso al arreglo en cada nivel, por lo que el tiempo continúa siendo $\Theta(n \log n)$.

\paragraph{Peor caso.} Ocurre cuando las mezclas deben alternar elementos de ambas mitades hasta casi agotarlas, con lo que se realizan aproximadamente $n\log_2 n - n + 1$ comparaciones (para $n$ potencia de 2). El tiempo es $\Theta(n \log n)$.

\paragraph{Caso promedio.} También es $\Theta(n \log n)$. Una característica importante de Merge Sort es que su tiempo de ejecución casi no depende del orden inicial de los datos.

\subsubsection{Complejidad en espacio}

La mezcla requiere una lista auxiliar de hasta $n$ posiciones para almacenar temporalmente los elementos ordenados (en la implementación, un arreglo \texttt{aux} del mismo tamaño que \texttt{arr}). Adicionalmente, la recursión utiliza una pila de llamadas de profundidad $O(\log n)$. En conjunto, la complejidad espacial auxiliar es
\[
  S(n) \in O(n).
\]
Por esta razón, Merge Sort \textbf{no} es un algoritmo in situ.

\subsection{Quick Sort (ordenamiento rápido)}
\subsubsection{Funcionamiento}

Quick Sort también sigue la estrategia \textbf{divide y vencerás}, pero a diferencia de Merge Sort el trabajo importante se hace al \emph{dividir} y no al combinar. Consta de tres etapas:

\begin{enumerate}
  \item \textbf{Dividir (particionar):} se elige un \textbf{pivote}, que en la implementación es el último elemento del subarreglo, \texttt{arr[hi]}. Luego se reorganiza \texttt{arr[lo..hi]} de modo que el pivote quede en su posición definitiva $p$, con todos los elementos de \texttt{arr[lo..p-1]} menores que el pivote y todos los de \texttt{arr[p+1..hi]} mayores o iguales.
  \item \textbf{Vencer:} las dos partes que quedaron a cada lado del pivote, \texttt{arr[lo..p-1]} y \texttt{arr[p+1..hi]}, se ordenan de forma recursiva con el mismo procedimiento. El \textbf{caso base} se da cuando \texttt{lo >= hi}, es decir, cuando el subarreglo tiene cero o un elemento, el cual ya está ordenado.
  \item \textbf{Combinar:} no hace falta ningún trabajo. Como la partición se hace sobre el propio arreglo y el pivote ya quedó en su lugar, cuando terminan las llamadas recursivas \texttt{arr[lo..hi]} ya está ordenado.
\end{enumerate}

La operación de partición (esquema de Lomuto) es el núcleo del algoritmo. Se utilizan dos índices: $j$ recorre el subarreglo desde \texttt{lo} hasta \texttt{hi - 1} e $i$ marca el final de la zona de elementos menores que el pivote (empieza en $\texttt{lo} - 1$, es decir, con la zona vacía).

\begin{enumerate}
  \item Se toma \texttt{pivote = arr[hi]} y se inicia \texttt{i = lo - 1}.
  \item Para cada $j$ de \texttt{lo} a \texttt{hi - 1} se compara \texttt{arr[j]} con el pivote. Si \texttt{arr[j] < pivote}, se incrementa $i$ y se intercambian \texttt{arr[i]} y \texttt{arr[j]}, de modo que el elemento entra a la zona de menores (si $i = j$ no hace falta mover nada). En caso contrario, \texttt{arr[j]} se queda donde está, en la zona de mayores o iguales.
  \item Al terminar el recorrido, se intercambia el pivote \texttt{arr[hi]} con \texttt{arr[i + 1]}. El pivote queda en la posición $p = i + 1$, que es la que ocupará en el arreglo ordenado.
\end{enumerate}

Durante el recorrido se cumple en todo momento que \texttt{arr[lo..i]} contiene elementos menores que el pivote, \texttt{arr[i+1..j-1]} contiene elementos mayores o iguales, y \texttt{arr[j..hi-1]} contiene los elementos que aún no se han revisado.

El procedimiento general consiste entonces en elegir un pivote, repartir los demás elementos a sus lados y repetir lo mismo con cada lado, que son subarreglos cada vez más pequeños, hasta que todos los elementos quedan en su posición definitiva. Esto es lo que muestra el visualizador: el pivote se marca en dorado, la zona de menores en verde azulado y la de mayores o iguales en violeta, con los índices $i$ y $j$ señalados debajo de las casillas y las posiciones definitivas en el color de ``ya ordenado''.

Debido a que la partición intercambia elementos que pueden estar muy separados entre sí, Quick Sort \textbf{no} es un algoritmo estable. Por ejemplo, con tres elementos iguales $1_a, 1_b, 1_c$ (en ese orden), esta implementación los deja como $1_c, 1_a, 1_b$.

\subsubsection{Complejidad en tiempo}

El costo de una partición sobre un subarreglo de $m$ elementos es $\Theta(m)$, pues cada elemento se compara con el pivote exactamente una vez ($m - 1$ comparaciones). Si el pivote queda con $k$ elementos a su izquierda, el tiempo satisface la recurrencia
\[
  T(n) = T(k) + T(n - k - 1) + \Theta(n), \qquad T(0) = T(1) = \Theta(1),
\]
por lo que el tiempo depende de qué tan equilibrada sea la división, es decir, de qué tan buena sea la elección del pivote.

\paragraph{Mejor caso.} Ocurre cuando el pivote cae siempre en medio del subarreglo, de modo que $k \approx n/2$ en cada partición. La recurrencia se convierte en $T(n) = 2\,T(n/2) + \Theta(n)$, cuya solución, por el Teorema Maestro, es $T(n) = \Theta(n \log n)$. Cuando $n = 2^k - 1$ y las divisiones son perfectas, se realizan exactamente $(n+1)\log_2(n+1) - 2n$ comparaciones.

\paragraph{Peor caso.} Ocurre cuando el pivote es siempre el menor o el mayor elemento del subarreglo: una de las partes queda vacía y la otra conserva $n - 1$ elementos. Entonces $T(n) = T(n-1) + \Theta(n)$ y se realizan $(n-1) + (n-2) + \cdots + 1 = n(n-1)/2$ comparaciones, es decir, $T(n) = \Theta(n^2)$. Con el último elemento como pivote esto sucede, por ejemplo, cuando el arreglo ya está ordenado (ascendente o descendente) o cuando todos los elementos son iguales (ninguno es menor que el pivote, así que todos van a la derecha).

\paragraph{Caso promedio.} Si todas las permutaciones de los datos son igualmente probables, el número esperado de comparaciones es
\[
  C(n) = 2(n+1)H_n - 4n,
\]
donde $H_n$ es el $n$-ésimo número armónico. Esta cantidad crece como $2n\ln n \approx 1.39\,n\log_2 n$, por lo que el tiempo promedio es $\Theta(n \log n)$. Incluso divisiones muy desiguales, como una proporción de 9 a 1, siguen produciendo $\Theta(n \log n)$: el comportamiento cuadrático solo aparece si las divisiones son extremas de forma sistemática.

A diferencia de Merge Sort, el tiempo de Quick Sort \textbf{sí} depende del orden inicial de los datos. En la práctica se reduce la probabilidad del peor caso eligiendo el pivote al azar o como la mediana de tres elementos; con pivote aleatorio, el tiempo esperado es $\Theta(n \log n)$ para cualquier entrada.

\subsubsection{Complejidad en espacio}

Quick Sort no necesita una lista auxiliar: la partición reorganiza los elementos dentro del propio arreglo mediante intercambios, por lo que cada llamada usa $O(1)$ de memoria adicional. Lo único que consume memoria es la pila de llamadas de la recursión, cuya profundidad es la altura del árbol de recursión: $O(\log n)$ en el mejor caso y en promedio, pero $O(n)$ en el peor caso, cuando las divisiones son extremas. Así,
\[
  S(n) \in O(\log n) \ \text{(caso promedio)}, \qquad S(n) \in O(n) \ \text{(peor caso)}.
\]
Por esta razón, Quick Sort se considera un algoritmo \textbf{in situ}, pues no usa un arreglo auxiliar, aunque no lo es en sentido estricto por la pila de recursión. Si se hace primero la llamada recursiva sobre la parte más pequeña y la otra parte se procesa con un ciclo, la profundidad de la pila queda acotada por $O(\log n)$ incluso en el peor caso (la implementación del visualizador no aplica esta optimización).

\subsection{Heap Sort (ordenamiento por montículo)}

\subsubsection{Montículo máximo}

Un montículo máximo, o \textit{max-heap}, es un árbol binario completo en el que cada padre es mayor o igual que sus hijos. Por esta propiedad, el valor máximo se encuentra en la raíz.

El árbol puede representarse en un arreglo. Si un nodo está en la posición \(i\), sus hijos se encuentran en las posiciones:

\[
\text{izquierdo}=2i+1, \qquad
\text{derecho}=2i+2
\]

Estas fórmulas usan índices que comienzan en cero. Heap Sort aprovecha esta representación para ordenar el mismo arreglo sin crear un árbol separado.

\subsubsection{Algoritmo}

El ordenamiento mediante montículos se realiza en dos etapas:

\begin{enumerate}
    \item Construir un montículo máximo a partir de la lista.
    \item Intercambiar la raíz con el último elemento pendiente, reducir el tamaño del montículo y restaurar su propiedad. Repetir hasta terminar.
\end{enumerate}

Para restaurar la propiedad del montículo, la función \texttt{heapify} compara una raíz con sus hijos, identifica el mayor y, si es necesario, realiza un intercambio. Después repite el procedimiento desde la posición del hijo intercambiado.

\begin{lstlisting}[caption={Restauración del montículo en Heap Sort}]
let mayor = raiz;
const izquierda = 2 * raiz + 1;
const derecha = 2 * raiz + 2;

if (izquierda < tamano) {
    if (a[izquierda] > a[mayor]) mayor = izquierda;
}
if (derecha < tamano) {
    if (a[derecha] > a[mayor]) mayor = derecha;
}
if (mayor !== raiz) {
    [a[raiz], a[mayor]] = [a[mayor], a[raiz]];
    heapify(tamano, mayor, ordenadas);
}
\end{lstlisting}

El siguiente fragmento muestra la construcción del montículo y la colocación sucesiva del máximo al final. Para simplificar, se omiten las llamadas que registran los pasos de la animación.

\begin{lstlisting}[caption={Construcción y extracción en Heap Sort}]
for (let raiz = Math.floor(a.length / 2) - 1;
     raiz >= 0; raiz--) {
    heapify(a.length, raiz, 0);
}

for (let final = a.length - 1; final > 0; final--) {
    [a[0], a[final]] = [a[final], a[0]];
    heapify(final, 0, a.length - final);
}
\end{lstlisting}

La construcción comienza en el último nodo que puede tener hijos y avanza hacia la raíz. En la segunda etapa, \texttt{final} determina el límite del montículo: los elementos que quedan a su derecha ya están ordenados.

\subsubsection{Objetivo y diseño de la animación}

El objetivo de esta animación es que se vea la relación entre el arreglo y el árbol que representa, y que se entienda cómo una raíz ``baja'' hasta recuperar la propiedad del montículo.

Para lograrlo, la página dibuja el arreglo como una fila de tarjetas numeradas y, además, el árbol de nodos con sus líneas padre--hijo. Los colores indican la raíz o inicio de la operación, el padre y el hijo que se están comparando, el mayor actual y la zona que ya quedó ordenada. Cuando un paso es un intercambio, el movimiento de los valores se anima para que el cambio sea visible.

\subsubsection{Análisis de complejidad}

Sea \(n\) la cantidad de elementos. Las siguientes complejidades corresponden al algoritmo de ordenamiento, sin incluir las copias del historial que utiliza la animación.

La construcción del montículo cuesta \(O(n)\). Después se realizan hasta \(n-1\) extracciones y cada restauración puede recorrer una altura de \(O(\log n)\). Por lo tanto, el tiempo total es:

\[
T(n)=O(n\log n)
\]

Este límite se mantiene en el mejor, promedio y peor caso. El ordenamiento modifica el arreglo en el mismo espacio; sin embargo, la versión utilizada llama recursivamente a \texttt{heapify}, por lo que puede ocupar \(O(\log n)\) de pila. Una versión iterativa podría usar \(O(1)\) de espacio auxiliar. Heap Sort no es estable en general.

\subsubsection{Ejemplo de ejecución}

Se parte del arreglo:

\[
[4,\,10,\,3,\,5,\,1]
\]

La siguiente tabla muestra el estado del arreglo después de cada intercambio registrado por \texttt{generar\-Pasos}. Los primeros dos intercambios construyen el montículo; los demás colocan el máximo pendiente al final y restauran el montículo.

\begin{table}[H]
\centering
\caption{Intercambios de Heap Sort para \([4,10,3,5,1]\)}
\begin{tabular}{cl}
\toprule
\textbf{Intercambio} & \textbf{Estado del arreglo} \\
\midrule
1 (construcción) & \([10,4,3,5,1]\) \\
2 (construcción) & \([10,5,3,4,1]\) \\
3 (extracción) & \([1,5,3,4,10]\) \\
4 (restauración) & \([5,1,3,4,10]\) \\
5 (restauración) & \([5,4,3,1,10]\) \\
6 (extracción) & \([1,4,3,5,10]\) \\
7 (restauración) & \([4,1,3,5,10]\) \\
8 (extracción) & \([3,1,4,5,10]\) \\
9 (extracción) & \([1,3,4,5,10]\) \\
\bottomrule
\end{tabular}
\end{table}

El resultado final es:

\[
\boxed{[1,\,3,\,4,\,5,\,10]}
\]

\subsubsection{Pasos registrados por la animación}

Para comprobar el funcionamiento de la lógica de la animación, se ejecutó la función \texttt{generar\-Pasos} con tres secuencias de 5, 10 y 15 elementos, y se contaron los pasos registrados y las operaciones que muestran los contadores de la página. En los tres casos, el último estado registrado fue el arreglo ordenado.

\begin{table}[H]
\centering
\caption{Pasos y operaciones registradas por la animación de Heap Sort}
\begin{tabular}{cccc}
\toprule
\(\boldsymbol{n}\) & \textbf{Pasos} & \textbf{Comparaciones} & \textbf{Intercambios} \\
\midrule
5 & 35 & 12 & 9 \\
10 & 112 & 41 & 30 \\
15 & 187 & 74 & 49 \\
\bottomrule
\end{tabular}
\end{table}

Heap Sort registra más pasos que Counting Sort porque cada llamada a \texttt{heapify} y cada comparación entre padre e hijo se guardan como un paso. Estos valores miden la cantidad de estados que genera la animación, no el tiempo de ejecución real del algoritmo.

\subsection{Comparación entre los algoritmos eficientes}

\begin{table}[htbp]
  \centering
  \caption{Comparación entre Merge Sort, Quick Sort y Heap Sort.}
  \begin{tabular}{lccc}
    \toprule
    \textbf{Característica} & \textbf{Merge} & \textbf{Quick} & \textbf{Heap} \\
    \midrule
    Mejor caso             & $\Theta(n \log n)$ & $\Theta(n \log n)$ & $\Theta(n \log n)$ \\
    Caso promedio          & $\Theta(n \log n)$ & $\Theta(n \log n)$ & $\Theta(n \log n)$ \\
    Peor caso              & $\Theta(n \log n)$ & $\Theta(n^2)$      & $\Theta(n \log n)$ \\
    Espacio aux. promedio  & $O(n)$             & $O(\log n)$        & $O(\log n)^{*}$ \\
    Espacio aux. peor caso & $O(n)$             & $O(n)$             & $O(\log n)^{*}$ \\
    Estable                & Sí                 & No                 & No \\
    In situ                & No                 & Sí (salvo pila)    & Sí (salvo pila)$^{*}$ \\
    Depende del orden      & No                 & Sí                 & Casi no \\
    Estrategia             & Divide y vencerás  & Divide y vencerás  & Montículo máximo \\
    \bottomrule
  \end{tabular}

  \smallskip
  {\footnotesize $^{*}$ En la implementación analizada, \emph{heapify} es recursivo; con una versión iterativa el espacio es $O(1)$.}
\end{table}

Los tres algoritmos alcanzan $\Theta(n \log n)$ en el caso promedio. Merge Sort y Heap Sort lo garantizan también en el peor caso: Merge Sort a cambio de memoria auxiliar $O(n)$ y Heap Sort sin arreglo auxiliar, pero sin ser estable. Quick Sort suele ser muy rápido en la práctica y casi no necesita memoria adicional, aunque su peor caso es $\Theta(n^2)$ (poco probable si se elige bien el pivote) y no es estable. Merge Sort es el único de los tres que es estable.

\clearpage

\section{Ordenamiento lineal}

\subsection{Counting Sort (ordenamiento por conteo)}

\subsubsection{Conteo de frecuencias}

Una frecuencia indica cuántas veces aparece un valor. Counting Sort almacena estas frecuencias en un arreglo auxiliar y después reconstruye la secuencia en orden.

Si los valores son enteros entre un mínimo y un máximo, se utiliza un desplazamiento para convertir cada valor en una posición del arreglo de conteo:

\[
\text{índice}=\text{valor}-\text{mínimo}
\]

De esta manera, también pueden representarse enteros negativos. La cantidad de posiciones necesarias para el conteo es:

\[
k=\text{máximo}-\text{mínimo}+1
\]

\subsubsection{Algoritmo}

Counting Sort no ordena comparando parejas de elementos. Primero encuentra el mínimo y el máximo, reserva un espacio para cada valor del rango y cuenta sus apariciones. Finalmente recorre las frecuencias desde el valor menor hasta el mayor y reconstruye el arreglo.

\begin{lstlisting}[caption={Conteo de apariciones en Counting Sort}]
const minimo = Math.min(...a);
const maximo = Math.max(...a);
const conteo = Array(maximo - minimo + 1).fill(0);

a.forEach((numero, indice) => {
    conteo[numero - minimo]++;
    foto(4, indice, numero, numero - minimo, 0);
});
\end{lstlisting}

El desplazamiento \texttt{numero - minimo} permite que el valor mínimo corresponda al índice cero. La función \texttt{foto} registra una copia del estado para mostrarlo posteriormente.

\begin{lstlisting}[caption={Reconstrucción del arreglo ordenado}]
let indiceResultado = 0;
for (let i = 0; i < conteo.length; i++) {
    const numero = i + minimo;
    for (let j = 0; j < conteo[i]; j++) {
        a[indiceResultado] = numero;
        indiceResultado++;
        foto(9, indiceResultado - 1, numero, i,
             indiceResultado);
    }
}
\end{lstlisting}

Cada valor se escribe tantas veces como indique su frecuencia. Esta versión reconstruye directamente valores numéricos; no conserva información adicional asociada a elementos iguales.

\subsubsection{Objetivo y diseño de la animación}

El objetivo es mostrar que este método no compara elementos entre sí: cada número se cuenta y después se escribe en la posición que le corresponde.

La vista tiene dos zonas. En la parte superior está el arreglo; en la inferior, una fila de tarjetas con un contador por cada valor del rango (por ejemplo, \texttt{x2} indica que ese número aparece dos veces). Durante la fase de conteo, una tarjeta naranja viaja del arreglo al contador de su valor y éste disminuye su cantidad restante en pantalla. En la fase de reconstrucción ocurre el recorrido inverso: la tarjeta sale del contador hacia su posición final y pasa a verde cuando queda ordenada.

\subsubsection{Análisis de complejidad}

Sea \(n\) la cantidad de elementos y \(k\) el tamaño del rango de valores enteros. Las siguientes complejidades corresponden al algoritmo de ordenamiento, sin incluir las copias del historial que utiliza la animación.

Se recorren los \(n\) elementos para encontrar el rango y contar sus frecuencias. Para reconstruir el resultado se recorren las \(k\) posiciones de conteo y se escriben los \(n\) valores:

\[
T(n,k)=O(n+k), \qquad S(n,k)=O(k)
\]

El espacio indicado corresponde al arreglo de frecuencias cuando se reconstruye sobre el arreglo de trabajo. Si el rango es mucho mayor que la cantidad de elementos, el método puede consumir memoria innecesaria. La variante que reconstruye solo números no conserva el orden de registros asociados a valores iguales; existen otras variantes de Counting Sort que sí pueden ser estables.

\subsubsection{Ejemplo de ejecución}

Se utiliza la lista:

\[
[4,\,2,\,2,\,1,\,3]
\]

El mínimo es 1, el máximo es 4 y el rango tiene cuatro posiciones. Las frecuencias son:

\begin{table}[H]
\centering
\caption{Frecuencias de los valores}
\begin{tabular}{ccccc}
\toprule
\textbf{Valor} & 1 & 2 & 3 & 4 \\
\midrule
\textbf{Frecuencia} & 1 & 2 & 1 & 1 \\
\bottomrule
\end{tabular}
\end{table}

Al escribir cada valor según su frecuencia se obtiene:

\[
\boxed{[1,\,2,\,2,\,3,\,4]}
\]

El número 2 aparece dos veces en el resultado, igual que en la entrada.

\subsubsection{Pasos registrados por la animación}

Igual que en Heap Sort, se ejecutó la función \texttt{generar\-Pasos} con secuencias de 5, 10 y 15 elementos. En los tres casos, el último estado registrado fue el arreglo ordenado.

\begin{table}[H]
\centering
\caption{Pasos y operaciones registradas por la animación de Counting Sort}
\begin{tabular}{cccc}
\toprule
\(\boldsymbol{n}\) & \textbf{Pasos} & \textbf{Comparaciones} & \textbf{Intercambios} \\
\midrule
5 & 14 & 0* & 0* \\
10 & 24 & 0* & 0* \\
15 & 34 & 0* & 0* \\
\bottomrule
\end{tabular}
\end{table}

\noindent *Counting Sort no compara ni intercambia elementos. La página usa esos contadores para mostrar los elementos contados (\(n\)) y los ya colocados.

Counting Sort registra \(2n+4\) pasos: la búsqueda del mínimo y del máximo, la creación del conteo, \(n\) conteos y \(n\) escrituras, más el estado final. El número de pasos crece linealmente con \(n\) en estas pruebas. Estos valores miden la cantidad de estados que genera la animación, no el tiempo de ejecución real del algoritmo.

\subsection{Radix Sort (ordenamiento por residuos)}
\subsubsection{Funcionamiento}

Radix Sort es un algoritmo que \textbf{no compara} los elementos entre sí: los ordena repartiéndolos en cubetas según sus dígitos. La versión que se estudia aquí es la LSD (\emph{least significant digit}) para números enteros no negativos en base $b = 10$: se procesa un dígito a la vez, empezando por el menos significativo (las unidades) y avanzando hacia el más significativo.

\begin{enumerate}
  \item \textbf{Preparar:} se determina el número de dígitos $d$ del mayor elemento $m$ del arreglo, $d = \lfloor \log_{10} m \rfloor + 1$. Ese será el número de pasadas.
  \item \textbf{Repetir:} se realizan $d$ pasadas, una por cada posición de dígito $p = 0, 1, \dots, d-1$ (unidades, decenas, centenas, \dots).
  \item \textbf{Terminar:} después de la última pasada el arreglo queda ordenado.
\end{enumerate}

Cada pasada utiliza 10 cubetas numeradas del 0 al 9, que son listas que respetan el orden de llegada. El dígito de un número $x$ en la posición $p$ es $\lfloor x / 10^{p} \rfloor \bmod 10$. La pasada tiene dos fases:

\begin{enumerate}
  \item \textbf{Repartir:} se recorre \texttt{arr} de izquierda a derecha y cada número se coloca al final de la cubeta que corresponde a su dígito en la posición $p$.
  \item \textbf{Recoger:} se recorren las cubetas de la 0 a la 9 y, dentro de cada una, los números en el orden en que llegaron; se van escribiendo de nuevo en \texttt{arr}, de izquierda a derecha.
\end{enumerate}

El algoritmo funciona porque, después de la pasada $p$, el arreglo queda ordenado según los $p + 1$ dígitos menos significativos de cada número. Esto se cumple porque cada pasada es \textbf{estable}: dos números con el mismo dígito en la posición actual conservan el orden relativo que traían de las pasadas anteriores, que ya los había ordenado por los dígitos de menor peso. Al terminar la pasada $d - 1$ se han considerado todos los dígitos y el arreglo está completamente ordenado.

Esto es lo que muestra el visualizador: la lista en la parte superior con el dígito analizado resaltado (los ceros a la izquierda se muestran atenuados), las diez cubetas debajo, una flecha que lleva cada número a su cubeta durante el reparto y, después, la recogida de las cubetas en orden.

Como las cubetas respetan el orden de llegada, Radix Sort es un algoritmo \textbf{estable}. Además, al no comparar elementos entre sí, no está sujeto a la cota inferior de $\Omega(n \log n)$ que tienen los algoritmos basados en comparaciones; a cambio, solo puede aplicarse a claves que se puedan descomponer en dígitos (enteros, cadenas de longitud acotada, etc.), y los enteros negativos requieren un tratamiento adicional.

\subsubsection{Complejidad en tiempo}

En cada pasada se recorren los $n$ elementos para repartirlos, lo que cuesta $\Theta(n)$, y después se recorren las $b$ cubetas para recogerlos, lo que cuesta $\Theta(n + b)$. Una pasada cuesta entonces $\Theta(n + b)$ y, como se realizan $d$ pasadas, el tiempo es
\[
  T(n) = \Theta\bigl(d\,(n + b)\bigr),
\]
con $b = 10$ y $d = \lfloor \log_{10} m \rfloor + 1$ dígitos.

\paragraph{Mejor caso.} Es $\Theta(d\,(n + b))$. El algoritmo siempre realiza las $d$ pasadas completas, de modo que incluso un arreglo que ya está ordenado se reparte y se recoge por completo en cada pasada.

\paragraph{Peor caso.} También es $\Theta(d\,(n + b))$. El tiempo no depende de cómo estén ordenados los datos, sino de $n$ y de $d$, es decir, de qué tan grandes son los números.

\paragraph{Caso promedio.} Igualmente $\Theta(d\,(n + b))$. Si el número de dígitos $d$ está acotado por una constante (por ejemplo, enteros de 32 bits), Radix Sort es lineal, $\Theta(n)$. Si, en cambio, los valores crecen con $n$ (por ejemplo, hasta $n^c$ con $c$ constante), entonces $d = \Theta(\log n)$ y el tiempo es $\Theta(n \log n)$. En general, Radix Sort supera a los algoritmos basados en comparaciones cuando $d$ es pequeño frente a $\log n$.

\subsubsection{Complejidad en espacio}

Durante el reparto, los $n$ elementos se encuentran dentro de las cubetas, además de las $b$ cubetas mismas, por lo que se necesita memoria auxiliar proporcional a $n + b$ (en la implementación, las diez cubetas juntas guardan los $n$ números de la pasada). No hay pila de recursión, pues el algoritmo es iterativo. La complejidad espacial auxiliar es
\[
  S(n) \in O(n + b) = O(n),
\]
ya que $b = 10$ es una constante. Por esta razón, Radix Sort \textbf{no} es un algoritmo in situ.

\subsection{Comparación entre los algoritmos lineales}

\begin{table}[htbp]
  \centering
  \caption{Comparación entre Counting Sort y Radix Sort.}
  \begin{tabular}{lcc}
    \toprule
    \textbf{Característica} & \textbf{Counting} & \textbf{Radix} \\
    \midrule
    Mejor caso          & $\Theta(n + k)$       & $\Theta(d\,(n + b))$ \\
    Caso promedio       & $\Theta(n + k)$       & $\Theta(d\,(n + b))$ \\
    Peor caso           & $\Theta(n + k)$       & $\Theta(d\,(n + b))$ \\
    Espacio auxiliar    & $O(k)$                & $O(n + b)$ \\
    Estable             & No (hay variantes estables) & Sí \\
    In situ             & No                    & No \\
    Compara elementos   & No                    & No \\
    Depende del orden   & No                    & No \\
    Tipo de datos       & Enteros en un rango acotado & Enteros (por dígitos) \\
    \bottomrule
  \end{tabular}

  \smallskip
  {\footnotesize $n$: número de elementos; $k$: tamaño del rango de valores; $d$: número de dígitos del mayor elemento; $b$: base (cubetas).}
\end{table}

Ninguno de los dos algoritmos compara elementos, por lo que no están sujetos a la cota $\Omega(n \log n)$ de los algoritmos basados en comparaciones. Counting Sort es lineal cuando el rango de valores $k$ es comparable con $n$, pero necesita un arreglo de conteo de tamaño $k$, por lo que su tiempo y su memoria crecen si el rango es muy grande. Radix Sort divide los números en dígitos y realiza $d$ pasadas estables, por lo que es lineal cuando $d$ es pequeño, a cambio de memoria adicional proporcional a $n$. Ambos solo sirven para claves enteras (o que se puedan descomponer en dígitos).

\clearpage

\section{Comparación de complejidades}

\begin{table}[htbp]
  \centering
  \caption{Complejidades en tiempo (mejor, promedio y peor caso) y en espacio auxiliar de los algoritmos de ordenamiento.}
  \begin{tabular}{lcccc}
    \toprule
    \textbf{Algoritmo} & \textbf{Mejor caso} & \textbf{Caso promedio} & \textbf{Peor caso} & \textbf{Espacio auxiliar} \\
    \midrule
    Bubble Sort    & $\Theta(n^2)^{*}$     & $\Theta(n^2)$         & $\Theta(n^2)$         & $O(1)$ \\
    Insertion Sort & $\Theta(n)$           & $\Theta(n^2)$         & $\Theta(n^2)$         & $O(1)$ \\
    Selection Sort & $\Theta(n^2)$         & $\Theta(n^2)$         & $\Theta(n^2)$         & $O(1)$ \\
    \midrule
    Merge Sort     & $\Theta(n \log n)$    & $\Theta(n \log n)$    & $\Theta(n \log n)$    & $O(n)$ \\
    Quick Sort     & $\Theta(n \log n)$    & $\Theta(n \log n)$    & $\Theta(n^2)$         & $O(\log n)$ prom.; $O(n)$ peor \\
    Heap Sort      & $\Theta(n \log n)$    & $\Theta(n \log n)$    & $\Theta(n \log n)$    & $O(\log n)^{\dagger}$ \\
    \midrule
    Counting Sort  & $\Theta(n + k)$       & $\Theta(n + k)$       & $\Theta(n + k)$       & $O(k)$ \\
    Radix Sort     & $\Theta(d\,(n + b))$  & $\Theta(d\,(n + b))$  & $\Theta(d\,(n + b))$  & $O(n + b)$ \\
    \bottomrule
  \end{tabular}

  \smallskip
  {\footnotesize $^{*}$ $\Theta(n)$ si se añade una bandera que detenga el algoritmo cuando una pasada no hace intercambios. \\
  $^{\dagger}$ Por la recursión de \emph{heapify} en la implementación analizada; con una versión iterativa es $O(1)$. \\
  $n$: número de elementos; $k$: tamaño del rango de valores; $d$: número de dígitos del mayor elemento; $b$: base (cubetas).}
\end{table}

\clearpage

\section{Problemas de LeetCode}

%------------------------------------------------------------
\subsection{LC 148: Ordenar una lista ligada}

\textbf{Problema.} Dada la cabeza de una lista ligada, ordenarla de forma ascendente en tiempo \bigO{n \log n}.

\textbf{Algoritmo.} La función \texttt{lc148} implementa \emph{merge sort} sobre listas. Si la lista está vacía o tiene un solo nodo, ya está ordenada. En caso contrario, \texttt{getMid} localiza el nodo central con dos apuntadores (uno avanza un nodo y el otro dos); la lista se corta en ese punto poniendo \texttt{der.next = None}, se ordenan recursivamente ambas mitades y finalmente \texttt{merge} las combina. \texttt{merge} utiliza un nodo ficticio (\texttt{dummy}) y recorre ambas listas enlazando siempre el menor de los dos nodos actuales; al terminar, anexa lo que sobre de la lista que no se agotó.

\textbf{Por qué funciona.} Cada nivel de recursión divide la lista a la mitad, por lo que hay \bigO{\log n} niveles, y en cada nivel la mezcla recorre todos los nodos una sola vez, es decir \bigO{n}; el total es \bigO{n \log n}. En \texttt{getMid}, \texttt{fast} inicia en \texttt{head.next}, de modo que \texttt{slow} termina en el final de la primera mitad y ambas mitades nunca quedan vacías. Como la mezcla solo reasigna apuntadores y no crea nodos, no se necesita memoria adicional más allá de la pila de recursión, \bigO{\log n}. La corrección se sigue por inducción: si las dos mitades están ordenadas, la mezcla de dos listas ordenadas produce una lista ordenada.

%------------------------------------------------------------
\subsection{LC 33: Búsqueda en un arreglo ordenado y rotado}

\textbf{Problema.} Se tiene un arreglo de enteros distintos, ordenado de forma ascendente pero rotado en algún pivote desconocido (por ejemplo, $[4,5,6,7,0,1,2]$). Dado un objetivo, devolver su índice o $-1$ si no está, en tiempo \bigO{\log n}.

\textbf{Algoritmo.} Es una búsqueda binaria modificada. En cada iteración se calcula el centro y, si coincide con el objetivo, se devuelve. Si no, se determina cuál mitad está ordenada: si \texttt{nums[izq] <= nums[centro]}, la mitad izquierda es la ordenada; de lo contrario lo es la derecha. Dentro de la mitad ordenada se verifica si el objetivo cae en su rango; si es así se continúa buscando ahí y si no, se descarta esa mitad y se busca en la otra.

\textbf{Por qué funciona.} Aunque el arreglo completo esté rotado, al partirlo por el centro \emph{al menos una de las dos mitades siempre está ordenada}, pues el punto de rotación solo puede estar en una de ellas. En una mitad ordenada se puede saber con certeza, mediante dos comparaciones, si el objetivo está o no en ella. Así, cada paso descarta la mitad del espacio de búsqueda, lo que da \bigO{\log n}.

%------------------------------------------------------------
\subsection{LC 1011: Capacidad para enviar paquetes en D días}

\textbf{Problema.} Se tiene una lista de pesos de paquetes que deben enviarse \emph{en el orden dado} en un barco, y un plazo de $D$ días. Cada día se carga el barco sin superar su capacidad. Encontrar la capacidad mínima con la que se pueden enviar todos los paquetes en a lo más $D$ días.

\textbf{Algoritmo.} Se hace una búsqueda binaria sobre la \emph{respuesta}, es decir, sobre la capacidad. El límite inferior es el paquete más pesado (menos que eso no se podría cargar ese paquete) y el superior es la suma de todos (todo en un día). Para una capacidad candidata, la función auxiliar \texttt{calcular\_dias} simula la carga: se acumulan paquetes y, cuando el siguiente ya no cabe, se inicia un nuevo día. Si la capacidad candidata permite terminar en $D$ días o menos, se guarda como posible respuesta y se busca una menor; si no, se busca una mayor.

\textbf{Por qué funciona.} La condición ``se puede enviar todo en $D$ días con capacidad $c$'' es monótona: si es posible con $c$, también lo es con cualquier capacidad mayor. Esto permite aplicar búsqueda binaria sobre un espacio de capacidades que va de \texttt{max(paquetes)} a \texttt{sum(paquetes)}, encontrando la frontera entre las capacidades insuficientes y suficientes. La simulación es voraz y óptima, porque llenar cada día lo más posible nunca aumenta el número de días. La complejidad es \bigO{n \log(S)}, donde $S$ es la suma de los pesos.

%------------------------------------------------------------
\subsection{LC 98: Validar un árbol binario de búsqueda}

\textbf{Problema.} Determinar si un árbol binario es un árbol binario de búsqueda (BST) válido: para cada nodo, todos los valores de su subárbol izquierdo son menores y todos los del derecho son mayores.

\textbf{Algoritmo.} La función auxiliar \texttt{valid(nodo, izq, der)} recorre el árbol recursivamente llevando un intervalo abierto $(izq, der)$ en el que debe encontrarse el valor del nodo. Inicia con $(-\infty, \infty)$. Un nodo nulo es válido; si el valor del nodo queda fuera del intervalo, el árbol no es válido. Al descender a la izquierda, el límite superior se actualiza al valor del nodo actual; al descender a la derecha, se actualiza el límite inferior.

\textbf{Por qué funciona.} Comparar únicamente cada nodo con sus hijos directos es insuficiente, pues un nodo del subárbol derecho podría ser menor que un ancestro lejano. Al propagar los límites hacia abajo, cada nodo se verifica contra \emph{todos} sus ancestros a la vez, ya que los límites acumulan la restricción impuesta por cada uno. Cada nodo se visita una vez: \bigO{n} en tiempo y \bigO{h} en espacio, con $h$ la altura del árbol.

%------------------------------------------------------------
\subsection{LC 322: Cambio de monedas}

\textbf{Problema.} Dadas monedas de distintas denominaciones (con existencia ilimitada) y una cantidad objetivo, calcular el número mínimo de monedas necesarias para formarla, o $-1$ si es imposible.

\textbf{Algoritmo.} Se usa programación dinámica \emph{bottom-up}. El arreglo \texttt{dp[a]} guarda el mínimo de monedas para formar la cantidad $a$. Se inicializa con \texttt{cantidad + 1} (un valor que funciona como ``infinito'', pues nunca se necesitarían más monedas que la cantidad misma usando monedas de valor 1) y \texttt{dp[0] = 0}. Para cada cantidad $a$ de $1$ a la cantidad objetivo y cada moneda $c \le a$ se aplica
\[
dp[a] = \min\bigl(dp[a],\; 1 + dp[a-c]\bigr).
\]
Al final, si \texttt{dp[cantidad]} sigue valiendo \texttt{cantidad + 1}, no hubo forma de alcanzarla y se devuelve $-1$.

\textbf{Por qué funciona.} El problema tiene subestructura óptima: la mejor manera de formar $a$ termina con alguna moneda $c$, y el resto $a-c$ debe formarse también de manera óptima. Al calcular las cantidades en orden creciente, cuando se necesita \texttt{dp[a - c]} ya está calculado. Además, los subproblemas se repiten, por lo que guardarlos evita la explosión exponencial de la recursión simple. La complejidad es \bigO{n \cdot m} con $n$ la cantidad y $m$ el número de monedas.

%------------------------------------------------------------
\subsection{LC 70: Subir escaleras}

\textbf{Problema.} Se sube una escalera de $n$ escalones, avanzando 1 o 2 escalones por paso. Determinar de cuántas formas distintas se puede llegar a la cima.

\textbf{Algoritmo.} Sea $f(n)$ el número de formas de llegar al escalón $n$. Se llega al escalón $n$ desde el $n-1$ (con un paso de 1) o desde el $n-2$ (con un paso de 2), por lo que
\[
f(n) = f(n-1) + f(n-2), \qquad f(0)=f(1)=1.
\]
Esta es la sucesión de Fibonacci. El código solo conserva los dos últimos valores (\texttt{uno} y \texttt{dos}) y los actualiza $n-1$ veces.

\textbf{Por qué funciona.} Las formas de llegar al escalón $n$ se dividen en dos conjuntos disjuntos según cuál fue el último paso, y cada conjunto tiene exactamente tantos elementos como formas de llegar al escalón previo correspondiente. Como cada valor solo depende de los dos anteriores, no hace falta guardar un arreglo completo, lo que da \bigO{n} en tiempo y \bigO{1} en espacio.

%------------------------------------------------------------
\subsection{LC 1396: Diseñar un sistema de metro}

\textbf{Problema.} Diseñar una clase \texttt{UndergroundSystem} que registre las entradas (\texttt{checkIn}) y salidas (\texttt{checkOut}) de clientes en estaciones y pueda responder el tiempo promedio de viaje entre dos estaciones (\texttt{getAverageTime}).

\textbf{Algoritmo.} Se usan dos diccionarios. \texttt{checkInMap} asocia el identificador de cada cliente con su estación y hora de entrada. \texttt{totalMap} asocia cada ruta (par estación de origen y estación de destino) con una lista \texttt{[tiempo total, número de viajes]}. En \texttt{checkOut} se recupera la entrada del cliente, se calcula la duración del viaje y se acumula en la ruta correspondiente. \texttt{getAverageTime} simplemente divide el tiempo total entre el número de viajes de esa ruta.

\textbf{Por qué funciona.} Para calcular un promedio no hace falta guardar cada viaje individual, basta con la suma y la cantidad. Así, las tres operaciones son \bigO{1} gracias al acceso directo de las tablas hash. Como un cliente solo puede estar en un viaje a la vez, el identificador es una llave suficiente para emparejar entrada y salida.

%------------------------------------------------------------
\subsection{LC 743: Tiempo de retraso de una red}

\textbf{Problema.} Dada una red de $n$ nodos con aristas dirigidas y ponderadas (tiempos de viaje) y un nodo inicial $k$, calcular el tiempo mínimo para que una señal emitida desde $k$ llegue a todos los nodos, o $-1$ si algún nodo no puede alcanzarse.

\textbf{Algoritmo.} Es el algoritmo de Dijkstra. Se construye una lista de adyacencia y se usa un min-heap de pares (tiempo acumulado, nodo), iniciando con $(0,k)$. Se extrae siempre el nodo con menor tiempo acumulado; si ya fue visitado se ignora, si no, se marca como visitado, se actualiza el tiempo total $t$ con el máximo visto y se insertan sus vecinos no visitados con el tiempo acumulado más el peso de la arista. Al final, si se visitaron los $n$ nodos se devuelve $t$ y si no, $-1$.

\textbf{Por qué funciona.} Con pesos no negativos, la primera vez que un nodo sale del heap lo hace con su distancia mínima desde $k$: cualquier otro camino a ese nodo pasaría por nodos con tiempo acumulado mayor o igual. Por tanto, la señal llega a todos los nodos cuando llega al más lejano, es decir, el máximo de las distancias mínimas. Si el número de nodos visitados es menor que $n$, existe algún nodo inalcanzable. La complejidad es \bigO{E \log E}, con $E$ el número de aristas.

%------------------------------------------------------------
\subsection{LC 210: Orden de cursos (ordenamiento topológico)}

\textbf{Problema.} Dados $n$ cursos y una lista de prerrequisitos (para tomar el curso $a$ hay que haber tomado antes el $b$), devolver un orden en el que se puedan tomar todos los cursos, o un arreglo vacío si es imposible.

\textbf{Algoritmo.} Se modela como un grafo dirigido y se realiza un ordenamiento topológico mediante búsqueda en profundidad (DFS). \texttt{prereq} guarda para cada curso la lista de sus prerrequisitos. La función \texttt{dfs} usa dos conjuntos: \texttt{cycle}, con los nodos de la ruta actual de recursión, y \texttt{visit}, con los cursos ya procesados por completo. Si se encuentra un curso que ya está en \texttt{cycle}, existe un ciclo y se devuelve un arreglo vacío. En caso contrario, se visitan primero todos los prerrequisitos del curso y, hasta que todos terminaron, el curso se agrega a \texttt{output}.

\textbf{Por qué funciona.} Un curso se agrega al resultado solamente después de agregar todos sus prerrequisitos (orden posterior del DFS), por lo que el resultado cumple por construcción que cada curso aparece después de sus requisitos. Un ordenamiento así existe si y solo si el grafo es acíclico, y un ciclo se manifiesta justamente como una arista de regreso hacia un nodo que aún está en la ruta de recursión, lo que detecta \texttt{cycle}. El conjunto \texttt{visit} evita reprocesar nodos, por lo que la complejidad es \bigO{V + E}.

%------------------------------------------------------------
\subsection{LC 198: Robo de casas}

\textbf{Problema.} Dada una lista con el dinero de cada casa, calcular la máxima cantidad que se puede robar sin tomar de dos casas adyacentes.

\textbf{Algoritmo.} Programación dinámica con dos variables: \texttt{anterior}, el mejor acumulado hasta la casa $i-2$, y \texttt{actual}, el mejor acumulado hasta la casa $i-1$. Para cada casa se decide entre saltarla (conservar \texttt{actual}) o robarla (sumar su valor a \texttt{anterior}):
\[
dp[i] = \max\bigl(dp[i-1],\; dp[i-2] + casa_i\bigr).
\]

\textbf{Por qué funciona.} Si se roba la casa $i$, no se puede haber robado la $i-1$, así que lo mejor que se puede combinar es el óptimo hasta la casa $i-2$; si no se roba, el óptimo es simplemente el de la casa $i-1$. Tomar el máximo de ambas opciones cubre todos los casos posibles. Como cada valor solo depende de los dos anteriores, basta con dos variables: \bigO{n} en tiempo y \bigO{1} en espacio.

%------------------------------------------------------------
\subsection{LC 218: Perfil de una ciudad (\emph{skyline})}

\textbf{Problema.} Dados los edificios como $[L, R, H]$ (coordenada izquierda, derecha y altura), devolver el contorno que forman en conjunto, como una lista de puntos clave $[x, altura]$ donde cambia la altura máxima visible.

\textbf{Algoritmo.} Es un barrido de línea con un max-heap. Cada edificio genera dos eventos: uno de inicio $(L, -H, R)$ y uno de fin $(R, 0, 0)$. Los eventos se ordenan por coordenada $x$; en caso de empate, los inicios con mayor altura van primero (por usar $-H$) y los eventos de fin después. El heap (implementado como min-heap con alturas negativas) guarda pares $(-H, R)$ de los edificios activos y comienza con un elemento centinela de altura 0 que nunca expira. En cada evento se eliminan de la cima los edificios cuyo borde derecho ya quedó atrás ($R \le x$), se agrega el edificio nuevo si es un inicio y se consulta la altura máxima actual. Si esta difiere de la última altura registrada, se agrega el punto $[x, altura]$ al resultado.

\textbf{Por qué funciona.} El contorno solo cambia en las coordenadas donde un edificio empieza o termina, y en cada una de ellas la altura visible es la del edificio activo más alto, que el heap entrega en la cima. La eliminación es \emph{perezosa}: un edificio terminado solo se retira cuando llega a la cima, ya que si está debajo de otro más alto no afecta la altura máxima. Registrar un punto únicamente cuando la altura cambia evita puntos redundantes en tramos consecutivos de igual altura. La complejidad es \bigO{n \log n}.

%------------------------------------------------------------
\subsection{LC 617: Combinar dos árboles binarios}

\textbf{Problema.} Dados dos árboles binarios, combinarlos sumando los valores de los nodos que se superponen; si en una posición solo hay un nodo, ese se conserva (el nulo cuenta como 0). Se debe devolver la raíz del nuevo árbol.

\textbf{Algoritmo.} Recursión sobre ambos árboles a la vez. Si ambos nodos son nulos, se devuelve \texttt{None}. En caso contrario, se toma el valor de cada nodo (0 si es nulo), se crea un nodo nuevo con la suma y sus hijos izquierdo y derecho se obtienen llamando recursivamente a la función con los hijos correspondientes de cada árbol (o \texttt{None} si ese nodo no existe).

\textbf{Por qué funciona.} La combinación de dos árboles se descompone de forma natural en combinar la raíz y combinar recursivamente los subárboles izquierdos y derechos entre sí. Tratar un nodo ausente como un valor 0 con hijos ausentes hace que el caso en que solo uno de los árboles continúa se copie sin tratamientos especiales. Cada posición se visita una vez: \bigO{\min(n_1, n_2)} en el caso de recorrido conjunto y, en el peor caso, \bigO{n_1 + n_2} al copiar ramas de un solo árbol. El código crea nodos nuevos, por lo que no modifica los árboles originales.

%------------------------------------------------------------
\subsection{LC 236: Ancestro común más bajo en un árbol binario}

\textbf{Problema.} Dado un árbol binario y dos nodos $p$ y $q$, encontrar su ancestro común más bajo (LCA), es decir, el nodo más profundo que tiene a ambos como descendientes (un nodo puede ser descendiente de sí mismo).

\textbf{Algoritmo.} Recursión posterior. Si el nodo actual es nulo o es $p$ o $q$, se devuelve tal cual. En otro caso se busca en el subárbol izquierdo y en el derecho. Si ambas búsquedas devuelven un nodo (no nulo), $p$ y $q$ están en lados distintos y la raíz actual es el LCA. Si solo una devuelve algo, se propaga hacia arriba ese resultado.

\textbf{Por qué funciona.} Cada llamada devuelve el primer nodo objetivo encontrado en su subárbol, o \texttt{None} si no hay ninguno. Cuando una llamada recibe resultados de ambos lados, esa es la primera posición, desde abajo, en la que ambos nodos quedan bajo el mismo ancestro, o sea el LCA. Si solo un lado devuelve algo, o bien se encontró uno solo de los nodos, o bien ambos están en ese lado y el resultado ya es el LCA, y propagarlo es correcto. El caso en que un nodo es ancestro del otro queda cubierto porque al encontrarlo se devuelve de inmediato, sin necesidad de buscar al otro. Se recorre cada nodo una vez: \bigO{n}.

%------------------------------------------------------------
\subsection{LC 235: Ancestro común más bajo en un BST}

\textbf{Problema.} Igual que el LC 236, pero el árbol es un árbol binario de búsqueda.

\textbf{Algoritmo.} Se aprovecha el orden del BST con un recorrido iterativo desde la raíz. Si los valores de $p$ y $q$ son ambos mayores que el del nodo actual, el LCA está en el subárbol derecho; si ambos son menores, está en el izquierdo. En cualquier otro caso (están en lados distintos, o uno de ellos es el nodo actual), el nodo actual es el LCA y se devuelve.

\textbf{Por qué funciona.} En un BST, todos los valores a la izquierda de un nodo son menores y los de la derecha mayores. Mientras $p$ y $q$ caigan del mismo lado, ambos siguen estando en ese subárbol y el LCA no puede estar fuera de él. El primer nodo en el que se ``separan'' es por definición su ancestro común más bajo. Como solo se desciende por una rama, la complejidad es \bigO{h} en tiempo y \bigO{1} en espacio, lo cual es mejor que el \bigO{n} del caso general.

%------------------------------------------------------------
\subsection{LC 200: Número de islas}

\textbf{Problema.} Dada una matriz de \texttt{"1"} (tierra) y \texttt{"0"} (agua), contar las islas, entendidas como grupos de celdas de tierra conectadas horizontal o verticalmente y rodeadas de agua.

\textbf{Algoritmo.} Se recorre toda la matriz. Cada vez que se encuentra una celda de tierra que no ha sido visitada, se incrementa el contador de islas y se lanza una búsqueda en anchura (BFS) desde ella. El BFS usa una cola y revisa las cuatro direcciones, agregando cada celda de tierra vecina no visitada y marcándola en el conjunto \texttt{visitados}.

\textbf{Por qué funciona.} El BFS desde una celda de tierra visita \emph{todas} las celdas conectadas a ella, es decir, la isla completa, y las marca. Gracias a esto, cuando el recorrido principal vuelve a encontrar una de esas celdas ya no la cuenta de nuevo. Por tanto, cada isla se cuenta exactamente una vez, cuando se encuentra su primera celda. Cada celda se procesa a lo más una vez, por lo que la complejidad es \bigO{f \cdot c}, con $f$ filas y $c$ columnas.

%------------------------------------------------------------
\subsection{LC 238: Producto del arreglo excepto uno mismo}

\textbf{Problema.} Dado un arreglo \texttt{nums}, devolver un arreglo donde la posición $i$ contenga el producto de todos los elementos excepto \texttt{nums[i]}, en tiempo lineal y sin usar división.

\textbf{Algoritmo.} Se hacen dos pasadas. En la primera, de izquierda a derecha, se guarda en \texttt{prods[i]} el producto de todos los elementos a la izquierda de $i$ (el prefijo), acumulándolo en una variable. En la segunda, de derecha a izquierda, se multiplica cada \texttt{prods[i]} por el producto de todos los elementos a su derecha (el posfijo), acumulado también en una variable.

\textbf{Por qué funciona.} El producto de todos los elementos excepto $nums[i]$ es
\[
\underbrace{\prod_{j<i} nums[j]}_{\text{prefijo}} \cdot \underbrace{\prod_{j>i} nums[j]}_{\text{posfijo}}.
\]
Cada pasada calcula una de las dos partes de forma incremental, de modo que no se necesita dividir el producto total entre $nums[i]$ (lo que además fallaría si hubiera ceros). Al reutilizar el arreglo de salida para el prefijo y solo una variable para el posfijo, el espacio adicional es \bigO{1} (sin contar la salida) y el tiempo es \bigO{n}.

%------------------------------------------------------------
\subsection{LC 253: Salas de reuniones}

\textbf{Problema.} Dada una lista de intervalos $[inicio, fin]$ de reuniones, calcular el número mínimo de salas necesarias para que todas se puedan realizar.

\textbf{Algoritmo.} Se separan los tiempos de inicio y de fin en dos arreglos que se ordenan de manera independiente. Con dos apuntadores, \texttt{s} sobre los inicios y \texttt{e} sobre los finales, se compara el siguiente inicio con el siguiente final. Si la reunión que comienza lo hace antes de que termine la más próxima a terminar, se necesita una sala más (el contador aumenta y se avanza \texttt{s}); si no, una sala se libera (el contador disminuye y se avanza \texttt{e}). La respuesta es el máximo valor que alcanza el contador.

\textbf{Por qué funciona.} El número de salas necesarias es el máximo número de reuniones que coinciden en un mismo instante. Al procesar, en orden cronológico, los eventos de inicio (+1) y de fin ($-1$), el contador es justamente el número de reuniones en curso, y su máximo es la respuesta. No importa saber qué inicio se empareja con qué final, solo cuántas reuniones están activas en cada momento, por lo que ordenar ambos arreglos por separado es válido. Como se usa una comparación estricta (\texttt{<}), una reunión que empieza justo cuando otra termina puede reutilizar su sala. La complejidad está dominada por el ordenamiento: \bigO{n \log n}.

\end{document}